\begingroup
\setlength{\parindent}{0pt}
\setlength{\parskip}{0.40\baselineskip}
\setlength{\abovedisplayskip}{6.5pt plus 1.5pt minus 1.5pt}
\setlength{\belowdisplayskip}{6.5pt plus 1.5pt minus 1.5pt}
\setlength{\jot}{5pt}
\clubpenalty=10000
\widowpenalty=10000

\section*{How the Proof Works}

The three-dimensional Navier--Stokes regularity problem asks whether a smooth
flow can become singular in finite time. We begin with a smooth,
incompressible flow of finite energy and a fixed positive viscosity. Its
velocity and pressure satisfy
\[
 \underbrace{\partial_tu}_{\text{local acceleration}}
 +\underbrace{(u\!\cdot\nabla)u}_{\text{nonlinear transport}}
 =\underbrace{-\nabla p}_{\text{pressure}}
 +\underbrace{\nu\Delta u}_{\text{viscous diffusion}},
 \qquad \nabla\!\cdot u=0.
\]
We start with a smooth, incompressible flow of finite energy and keep its
viscosity fixed and positive. We know it stays smooth for a short time; the
question is whether it stays smooth forever or becomes singular in finite
time. The hope is simple: viscosity smooths the fluid faster than nonlinear
transport can sharpen it.

Consider the case of a vortex within the fluid. Taking the curl of momentum
gives us:
\[
 D_t\omega=(\omega\cdot\nabla)u+\nu\Delta\omega,
 \qquad \omega=\nabla\times u,
\]
where \(D_t=\partial_t+u\cdot\nabla\).

Strain within the fluid's flow can naturally stretch a vortex. When that
happens, incompressibility forces a material tube to narrow as it lengthens.
If circulation persists strongly enough through that narrowing, the fluid
circulates faster around a smaller core.

For a material tube that stays locally cylindrical, let \(\ell\) and \(r\)
be its length and radius, \(\Gamma\) its circulation,
\(\sigma=\dot\ell/\ell>0\) its axial stretching rate, and
\(v=|\Gamma|/(2\pi r)\) its mean swirl speed. Following the fluid,
\[
 \begin{aligned}
 \frac{d}{dt}(\pi r^2\ell)=0
 &\quad\Longrightarrow\quad \frac{\dot r}{r}=-\frac{\sigma}{2},\\
 \frac{\dot v}{v}
 &=\frac{d}{dt}\log|\Gamma|+\frac{\sigma}{2}>0
 \quad\text{if}\quad \frac{d}{dt}\log|\Gamma|>-\frac{\sigma}{2}.
 \end{aligned}
\]

For a material tube following a chosen group of fluid particles: a vortex
core is the region where rotation is concentrated. Stretching squeezes
those particles inward. Viscosity diffuses momentum between neighbouring
regions of fluid, reducing velocity differences and spreading that rotation
through the surrounding fluid.
The material tube can narrow while some rotation spreads beyond its
boundary, reducing the circulation enclosed by it. So narrowing the vortex
core makes the fluid circulate faster only if the enclosed circulation does
not weaken too quickly.

Whether viscosity can smooth those differences fast enough is the difficulty.

A narrowing vortex makes that difficulty progressively more local. Large
gradients can occupy a region too small to account for much of the total
energy, while both nonlinear transport and viscous diffusion become stronger
there. Following the concentration toward vanishing width forces us to
compare their behavior at smaller and smaller scales. We need to know whether viscosity gains on transport as the width tends
to zero.

The equations allow an exact comparison. For a scale factor \(\lambda>1\),
shrink lengths by \(\lambda\), increase velocities by \(\lambda\), and
shorten times by \(\lambda^2\):
\[
 u_\lambda(x,t)=\lambda u(\lambda x,\lambda^2t),
 \qquad
 p_\lambda(x,t)=\lambda^2p(\lambda x,\lambda^2t).
\]
The time derivative, nonlinear transport, pressure gradient, and viscous
term all acquire a factor \(\lambda^3\). After that common factor is
removed, the viscosity is still \(\nu\). The narrower, faster flow obeys
the original equation with precisely the original balance of scaling
weights. Shrinking the scale has supplied no extra viscous advantage.
Which effect dominates the evolution still depends on the fluid's actual
configuration.

To measure what this concentration does to spatial regularity, the Sobolev
index \(s\) lets us vary how strongly we count fine spatial variation.
At integer values it counts square-integrable derivatives; fractional values
interpolate that sensitivity by weighting a spatial frequency \(\xi\) by
\(\lvert\xi\rvert^s\). The velocity factor contributes \(\lambda\), the
derivative weight contributes \(\lambda^s\), and the square root of the
volume factor contributes \(\lambda^{-3/2}\). Together they give
\[
 \|u_\lambda(\cdot,t)\|_{\dot H^s}
 =\lambda^{s-1/2}\|u(\cdot,\lambda^2t)\|_{\dot H^s}.
\]
Below \(s=\tfrac12\), concentration makes this norm smaller. Above
\(s=\tfrac12\), it makes the norm larger. Exactly at \(s=\tfrac12\),
the increased sensitivity to fine variation offsets the shrinking volume,
and the norm is unchanged. This is the critical height. Its invariance
compares rescaled flows; it is not a conservation law along an evolving flow.

Energy lies below that height. It is half the squared \(L^2\) norm, so it
scales as \(\lambda^{-1}\): the velocity grows, but the volume containing the fast
motion shrinks even faster. Increasing concentration and speed can coexist
with falling total energy. The energy identity can keep the total finite
without, by that fact alone, excluding an arbitrarily narrow and fast region.

The time available for further concentration also shrinks. At width
\(\ell\) and velocity scale \(U\), transport acts over a time of order
\(\ell/U\); under this rescaling that time gains a factor
\(\lambda^{-2}\). The viscous time \(\ell^2/\nu\) gains that factor too.
If the evolution could repeat such stages, each narrower stage would take
less time than the last:
\[
 \ell_j=\lambda^{-j}\ell_0,
 \qquad U_j=\lambda^jU_0,
 \qquad \tau_j=\lambda^{-2j}\tau_0,
 \qquad
 \sum_{j=0}^{\infty}\tau_j
 =\frac{\tau_0}{1-\lambda^{-2}}<\infty.
\]
Infinitely many stages could then fit before a finite terminal time \(T_*\). That
is the hypothetical Zeno cascade: narrowing without a smallest width,
acceleration without a final stage, and finite energy throughout. Scaling
allows this comparison. It does not establish that a Navier--Stokes flow
actually performs the cascade.

For a finite-time singularity, velocity would have to blow up. That requires
acceleration to blow up, then jerk, and so on, and so on. Here we follow the
velocity field at fixed spatial positions, so the sequence is
\(u,\partial_tu,\partial_t^2u,\ldots\). These are Eulerian time derivatives;
a moving particle has material acceleration
\(D_tu=\partial_tu+(u\cdot\nabla)u\), and its later responses use
\(D_t=\partial_t+u\cdot\nabla\).
The necessity has a precise meaning in the spatial \(L^\infty\) norm.
For any fixed \(t_0<T_*\) and every finite order \(j\), integration gives
\[
 \|\partial_t^ju(t)\|_\infty
 \le \|\partial_t^ju(t_0)\|_\infty
 +\int_{t_0}^{t}\|\partial_t^{j+1}u(s)\|_\infty\,ds.
\]
A bounded next derivative could accumulate only a finite change before
\(T_*\). Unbounded velocity requires an infinite accumulation of its time
derivative, which requires that derivative to be unbounded; applying the
argument again forces the next order to be unbounded as well. The norms
need not grow monotonically, and their large values need not occur at one
fixed point. The requirement persists through every finite temporal order.

But these increasingly demanding time derivatives are still derivatives of
a velocity governed by Navier--Stokes. Each time we differentiate, the
viscous term applies a Laplacian to the preceding response. The first time
derivative contains \(\Delta u\), the second contains
\(\Delta\partial_tu\), and substitution of the first equation brings in
\(\Delta^2u\). Continuing exposes higher spatial derivatives, together
with all the products created by differentiating transport and pressure.
The temporal escalation has reached into the spatial structure of the flow.

We already measured the velocity's fine spatial variation with the critical
norm. Differentiating that norm follows how this spatial variation changes
in time, so it brings the temporal responses directly into the measurement.
With \(\Lambda=(-\Delta)^{1/2}\), write half its squared size as
\(H=\tfrac12\|\Lambda^{1/2}u\|_2^2\). Differentiating the two velocity
factors gives
\[
 \begin{aligned}
 H'&=\operatorname{Re}\langle\Lambda^{1/2}u,
                          \Lambda^{1/2}\partial_tu\rangle,\\
 H''&=\|\Lambda^{1/2}\partial_tu\|_2^2
 +\operatorname{Re}\langle\Lambda^{1/2}u,
                          \Lambda^{1/2}\partial_t^2u\rangle.
 \end{aligned}
\]
Acceleration and jerk enter through these actual differentiated velocity
factors. The equation now expresses those factors through spatial
derivatives, transport, and pressure. In the first pairing, its viscous
contribution is
\[
 \nu\langle\Lambda^{1/2}u,\Lambda^{1/2}\Delta u\rangle
 =-\nu\|\Lambda^{3/2}u\|_2^2.
\]
The Laplacian has converted the temporal response into a measurement of
finer spatial variation. Integration by parts does this at any Sobolev
height: if \(E_s=\tfrac12\|\Lambda^su\|_2^2\), then
\[
 E_s'=-2\nu E_{s+1}+\mathcal P_s,
 \qquad
 \mathcal P_s:=-\operatorname{Re}\langle\Lambda^su,
 \Lambda^s((u\cdot\nabla)u+\nabla p)\rangle_{L^2}.
\]
Differentiate the critical-height relation again. The new spatial energy
\(E_{3/2}\) changes according to this identity too, exposing
\(E_{5/2}\); its transport and pressure terms also have to be differentiated.
The first two orders become
\[
 \begin{aligned}
 H'&=-2\nu E_{3/2}+\mathcal P_{1/2},\\
 H''&=4\nu^2E_{5/2}-2\nu\mathcal P_{3/2}
       +\partial_t\mathcal P_{1/2}.
 \end{aligned}
\]
The two expressions for \(H''\) concern one quantity: the pairing of velocity
with its higher temporal responses also contains its higher spatial
Sobolev energy, with the differentiated nonlinear contribution retained.
Continuing repeats this relation at every finite order. The temporal
responses and spatial regularity cannot be specified independently of one
another; the equation couples them through these identities. Their bounds
must account for the nonlinear terms as well as the viscous term.

Those spatial demands are exactly what smoothness needs. A bound at one
Sobolev height controls only a finite amount of spatial variation. Raise
the height far enough, however, and Sobolev embedding makes any prescribed
spatial derivative continuous: to obtain \(k\) continuous derivatives,
choose \(m>k+\tfrac32\). If the velocity belongs to every finite height,
we can make that choice for every \(k\):
\[
 L^2\cap\bigcap_{n\ge0}\dot H^{n+1/2}
 =\bigcap_{m\ge0}H^m
 =H^\infty
 \hookrightarrow C^\infty.
\]
No single finite height gives all of smoothness. The intersection does,
because every derivative we ask for has a sufficiently high finite space
in which it is controlled. The constants may increase with the order;
smoothness asks that each finite order exists, with its own finite bound.

The equation also explains why this spatial intersection can accommodate
the entire temporal sequence. A Laplacian uses two spatial derivatives;
at any requested height, the intersection has those additional derivatives
available. Products of smooth Sobolev fields and the pressure recovered from
them also remain in the intersection. Once the spatial and temporal
requirements have been realized together over the approach to \(T_*\),
the equation determines the next temporal response there, and its
differentiated equations continue that relation through every finite order.
The time and space requirements meet in a velocity with enough spatial
regularity for every temporal response and enough time regularity to reach
a common endpoint. That endpoint is smooth and can serve as initial data
for continued evolution. This is the connection the proof has to realize
from the initial velocity \(u_0\); the differentiated identities reveal it, while the construction
must establish the bounds and compatibility that make it true at the endpoint.

If a singularity were to form, we would expect the opposite behavior on
the approach to that time. The velocity could remain smooth at every earlier
instant while its increasingly fine structure forced some fixed spatial
Sobolev norm beyond every bound. Its temporal responses would become
unbounded as well, even while the total energy remained finite. We would
have all finite derivatives before \(T_*\), yet lose the control needed to
carry them through \(T_*\). The distinction is between smoothness on each
shorter interval and regularity maintained over the full approach to the
putative singularity. A proof has to exclude this loss at the endpoint.

\needspace{11\baselineskip}
To do that, we return to the derivatives the initial velocity actually
generates. Pressure is fixed by incompressibility and its decay
normalization: taking the divergence of the equation gives
\(-\Delta p=\partial_i\partial_j(u_i u_j)\). Differentiating the velocity
equation now gives the first three responses in their uncompressed form:
\begin{flalign*}
 &\partial_tu=\nu\Delta u-(u\cdot\nabla)u-\nabla p,&&\\
 &\partial_t^2u=\nu\Delta\partial_tu
   -(\partial_tu\cdot\nabla)u-(u\cdot\nabla)\partial_tu
   -\nabla\partial_tp,&&\\
 &\partial_t^3u=\nu\Delta\partial_t^2u
   -(\partial_t^2u\cdot\nabla)u
   -2(\partial_tu\cdot\nabla)\partial_tu
   -(u\cdot\nabla)\partial_t^2u-\nabla\partial_t^2p.&&
\end{flalign*}
The factor of two comes from differentiating both factors of transport.
At the next order each product splits again. Writing
\(V_r=\partial_t^ru\) and \(Q_r=\partial_t^rp\) lets us retain all those
terms without expanding an ever longer line:
\[
 \begin{aligned}
 V_{r+1}&=\nu\Delta V_r
 -\sum_{a=0}^r\binom ra(V_a\cdot\nabla)V_{r-a}-\nabla Q_r,\\
 Q_r&=R_iR_j\sum_{a=0}^r\binom ra V_{a,i}V_{r-a,j}.
 \end{aligned}
\]
Here \(R_i\) are the Riesz transforms and repeated spatial indices are
summed. The pressure formula differentiates the incompressibility constraint
along with the velocity. At \(t=0\), these relations start with \(V_0=u_0\)
and determine each subsequent velocity and pressure derivative. Spatial
differentiation fills out each temporal row. This collection is the mixed
jet: the time and space derivatives generated by the datum under the equation.

We need only finitely many of those derivatives for any specified order of
regularity. Fix a temporal depth \(N\) and a spatial height \(M\). The
datum-generated rectangle theorem, Theorem~\ref{thm:datum-rectangles},
supplies the adjacent rows on the whole interval ending at \(T_*\).
In the cutoff form of its conclusion,
\[
 \sup_{0<T<T_*}\sum_{n=0}^{N}\left(
 \|V_n\|_{L^2(0,T;H^{M+1})}^2
 +\|V_{n+1}\|_{L^2(0,T;H^{M-1})}^2
 \right)<\infty.
\]
The bound can depend on \(N,M,u_0,\nu\), and the putative finite endpoint,
but it holds as the cutoff \(T\) approaches that endpoint. Increasing the
derivative orders may increase the bound; approaching \(T_*\) at any fixed
pair of orders cannot make it diverge. The differentiated recurrence
identifies every entry with a derivative of the original history, while the
datum theorem supplies its full-interval control.

Letting \(T\) increase to \(T_*\) gives
\(V_n\in L^2(0,T_*;H^{M+1})\) and
\(\partial_tV_n=V_{n+1}\in L^2(0,T_*;H^{M-1})\).
These adjacent spaces are chosen for a reason. The time derivative can be
paired with \(V_n\) across the intermediate space \(H^M\), and the two
square-integrable factors give an integrable change in its squared norm.
The Hilbert-triple trace theorem makes this precise:
\[
 H^{M+1}\subset H^M\subset H^{M-1},
 \qquad V_n\in C([0,T_*];H^M).
\]
Each retained velocity derivative now has an actual endpoint value at the
requested spatial height. The pressure formula supplies the corresponding
pressure rows. We collect these finitely many velocity, adjacent time,
and pressure rows into a rectangle \(\mathcal R_{N,M}[u_0;T_*]\).
Its two indices remain independent, since a request for another time
derivative and a request for finer spatial regularity are different requests.

\needspace{10\baselineskip}
Increasing either depth retains every derivative already present. An entry
in two rectangles is a distributional derivative of the original velocity
or pressure in both, so their restrictions agree. This matters at the
endpoint: different spatial heights must describe one limiting field,
and different temporal depths must describe its actual derivatives.
Compatibility lets all the finite rectangles be retained together:
\[
 \mathcal I[u_0]
 :=\bigl(\mathcal R_{N,M}[u_0;T_*]\bigr)_{N,M\ge0},
 \qquad
 u\in\bigcap_{r,m\ge0}H^r(0,T_*;H^m(\mathbb R^3)).
\]
For any finite \(r,m\), choose a rectangle deep enough to contain those
derivatives. Every larger rectangle preserves them. The projective
intersection expresses this agreement through all finite orders; it does
not require summing their bounds into one infinite norm.

We can now take the endpoint in every spatial height. The zeroth temporal
row gives \(u\in C([0,T_*];H^m)\) for every finite \(m\), and compatibility
identifies its traces as one field:
\[
 u_*:=\lim_{t\uparrow T_*}u(t)
 \quad\hbox{in every }H^m,
 \qquad u_*\in H^\infty_\sigma\hookrightarrow C^\infty.
\]
The differentiated velocity and pressure relations also reach this endpoint,
determining its full jet from \(u_*\). Smooth local existence then starts a
solution from \(u_*\). Local uniqueness on overlapping time intervals joins
that solution to the original flow approaching \(T_*\), continuing it beyond
its supposed maximal time. A finite \(T_*\) is impossible, so \(T_*=\infty\).

With the intersection constructed, we can derive the finite estimates,
and use them. Any estimate requiring finitely many derivatives can take its
regularity from a sufficiently large rectangle. For a finite time \(T\),
temporal order \(r\), spatial order \(k\), and \(2\le p\le\infty\), choose
\(m>k+\tfrac32-\tfrac3p\). The corresponding coordinate and Sobolev embedding
give
\[
 \partial_t^ru\in C([0,T];H^m)
 \hookrightarrow L^\infty(0,T;W^{k,p}).
\]
Taking \(r=0\), \(k=1\), \(p=\infty\), and \(m=3\), for example, gives
\[
 \int_0^T\|\nabla u(t)\|_\infty\,dt
 \le T\sup_{0\le t\le T}\|\nabla u(t)\|_\infty<\infty.
\]
The accumulated velocity gradient is finite on every finite interval.
This is the familiar control that a concentrating vortex would have to
defeat, now obtained from a finite projection of the completed intersection.

\endgroup
